%=====================================================================
%  We Know Group | Kato Kiosk Configuration Guide  (AM + IT)
%  Command Centre setup, silent app install, kiosk lockdown.
%  Shared styling lives in wkg-preamble.tex.
%=====================================================================
\documentclass[a4paper,10pt]{article}
%=====================================================================
%  We Know Group | Shared preamble for the Kato kiosk guides
%
%  Branding per "We Know Group Brand Guidelines V2 2023":
%    - Primary colours: Graphite #121212, Azure #00A6D3, White
%    - Secondary: Sky Blue #E5F1F9, Sea Blue #008BB2, Night Blue #0C2752
%    - Greys: Night #575A5C, Mid #ADADAD, Light #F4F4F4
%    - Typeface: Circular Pro (approximated here with Inter, the licensed
%      brand stand-in used in the WeKnow Employee Handbook)
%    - Logos are the official WeKnow assets in images/ -- never redrawn.
%
%  Both Kato kiosk guides share this file so their branding stays in step.
%
%  Usage, from a document that has already declared \documentclass:
%      %=====================================================================
%  We Know Group | Shared preamble for the Kato kiosk guides
%
%  Branding per "We Know Group Brand Guidelines V2 2023":
%    - Primary colours: Graphite #121212, Azure #00A6D3, White
%    - Secondary: Sky Blue #E5F1F9, Sea Blue #008BB2, Night Blue #0C2752
%    - Greys: Night #575A5C, Mid #ADADAD, Light #F4F4F4
%    - Typeface: Circular Pro (approximated here with Inter, the licensed
%      brand stand-in used in the WeKnow Employee Handbook)
%    - Logos are the official WeKnow assets in images/ -- never redrawn.
%
%  Both Kato kiosk guides share this file so their branding stays in step.
%
%  Usage, from a document that has already declared \documentclass:
%      %=====================================================================
%  We Know Group | Shared preamble for the Kato kiosk guides
%
%  Branding per "We Know Group Brand Guidelines V2 2023":
%    - Primary colours: Graphite #121212, Azure #00A6D3, White
%    - Secondary: Sky Blue #E5F1F9, Sea Blue #008BB2, Night Blue #0C2752
%    - Greys: Night #575A5C, Mid #ADADAD, Light #F4F4F4
%    - Typeface: Circular Pro (approximated here with Inter, the licensed
%      brand stand-in used in the WeKnow Employee Handbook)
%    - Logos are the official WeKnow assets in images/ -- never redrawn.
%
%  Both Kato kiosk guides share this file so their branding stays in step.
%
%  Usage, from a document that has already declared \documentclass:
%      \input{wkg-preamble}
%      \renewcommand{\guidetitle}{Short title for the page footer}
%=====================================================================

%---------------------------------------------------------------------
% Fonts & encoding -- Inter (brand stand-in for Circular Pro)
%---------------------------------------------------------------------
\usepackage[T1]{fontenc}
\usepackage[utf8]{inputenc}
\usepackage[sfdefault]{inter}
\renewcommand{\familydefault}{\sfdefault}
\usepackage{microtype}

%---------------------------------------------------------------------
% Layout
%---------------------------------------------------------------------
\usepackage[a4paper,margin=20mm,top=28mm,bottom=22mm,headheight=20pt,headsep=14pt]{geometry}
\usepackage{graphicx}
\graphicspath{{images/}}
\usepackage[table]{xcolor}
\usepackage{tikz}
\usetikzlibrary{arrows.meta,positioning,calc,shapes.misc,fit,backgrounds}
\usepackage{tcolorbox}
\tcbuselibrary{skins,breakable}
\usepackage{booktabs}
\usepackage{array}
\usepackage{tabularx}
\usepackage{enumitem}
\usepackage{fancyhdr}
\usepackage{titlesec}
\usepackage{listings}
\usepackage{float}
\usepackage{lastpage}
\usepackage{amssymb}
\usepackage{caption}
\usepackage[hidelinks,colorlinks=true,urlcolor=weknowsea,linkcolor=weknownight]{hyperref}

%---------------------------------------------------------------------
% We Know Group palette (Brand Guidelines V2 2023, Section 3)
%---------------------------------------------------------------------
\definecolor{weknowazure}{HTML}{00A6D3}     % Hero colour
\definecolor{weknowgraphite}{HTML}{121212}
\definecolor{weknowsky}{HTML}{E5F1F9}
\definecolor{weknowsea}{HTML}{008BB2}
\definecolor{weknownight}{HTML}{0C2752}
\definecolor{weknowgreyN}{HTML}{575A5C}
\definecolor{weknowgreyM}{HTML}{ADADAD}
\definecolor{weknowgreyL}{HTML}{F4F4F4}
\definecolor{weknowazure80}{HTML}{33B8DC}    % Azure tints (brand: tint Azure only)
\definecolor{weknowazure60}{HTML}{66CAE5}
\definecolor{weknowazure40}{HTML}{99DBED}
\definecolor{weknowazure20}{HTML}{CCEDF6}

%---------------------------------------------------------------------
% Headers / footers -- \guidetitle is set by each document
%---------------------------------------------------------------------
\providecommand{\guidetitle}{Kato Kiosk Guide}
\pagestyle{fancy}
\fancyhf{}
% \sectionmark uses \markright, so the running head must read \rightmark.
\fancyhead[L]{\footnotesize\color{weknowgreyN}\nouppercase{\rightmark}}
\fancyhead[R]{\raisebox{-0.35\height}{\includegraphics[height=15pt]{wkg-logo-horizontal}}}
\renewcommand{\headrulewidth}{0.8pt}
\renewcommand{\headrule}{\hbox to\headwidth{\color{weknowazure}\leaders\hrule height \headrulewidth\hfill}}
\fancyfoot[L]{\scriptsize\color{weknowgreyN}Internal --- We Know Group}
\fancyfoot[C]{\scriptsize\color{weknowgreyN}\guidetitle}
\fancyfoot[R]{\scriptsize\color{weknowgreyN}Page \thepage\ of \pageref{LastPage}}
\renewcommand{\footrulewidth}{0.4pt}
\renewcommand{\footrule}{\hbox to\headwidth{\color{weknowgreyM}\leaders\hrule height \footrulewidth\hfill}}

%---------------------------------------------------------------------
% Section styling
%---------------------------------------------------------------------
\titleformat{\section}
  {\color{weknowgraphite}\LARGE\bfseries}
  {\color{weknowazure}\thesection}{0.7em}{}
  [{\vspace{3pt}\color{weknowazure}\titlerule[1.4pt]}]
\titleformat{\subsection}
  {\color{weknowsea}\Large\bfseries}
  {\color{weknowazure}\thesubsection}{0.6em}{}
\titlespacing*{\section}{0pt}{20pt}{12pt}
\titlespacing*{\subsection}{0pt}{16pt}{6pt}
% The running head takes its text from \wksecmark, not from the section
% title, so role badges never leak into the header. \wksection and
% \wkrolesection (defined below) set it. Subsections must not overwrite
% the section in the head.
\newcommand{\wksecmark}{}
\renewcommand{\sectionmark}[1]{\markright{\thesection\quad\wksecmark}}
\renewcommand{\subsectionmark}[1]{}

%---------------------------------------------------------------------
% Captions
%---------------------------------------------------------------------
\captionsetup{font=small,labelfont={bf,color=weknownight},
              textfont={color=weknowgreyN},skip=5pt}

%---------------------------------------------------------------------
% Screenshot helper -- \screenshot[width fraction]{file}{caption}
%---------------------------------------------------------------------
\newcommand{\screenshot}[3][0.72]{%
  \begin{figure}[H]
    \centering
    \setlength{\fboxsep}{0pt}\setlength{\fboxrule}{0.6pt}%
    {\color{weknowgreyM}\fbox{\includegraphics[width=#1\linewidth]{#2}}}%
    \caption{#3}
  \end{figure}}

%---------------------------------------------------------------------
% Callout boxes
%---------------------------------------------------------------------
\newtcolorbox{wkinfo}[1][]{%
  enhanced,breakable,colback=weknowsky,colframe=weknowazure,
  boxrule=0pt,leftrule=2.5pt,arc=1.5pt,
  fonttitle=\bfseries,coltitle=weknownight,title={#1},
  left=6pt,right=6pt,top=5pt,bottom=5pt}
\newtcolorbox{wkwarn}[1][]{%
  enhanced,breakable,colback=weknowgreyL,colframe=weknownight,
  boxrule=0pt,leftrule=2.5pt,arc=1.5pt,
  fonttitle=\bfseries,coltitle=white,title={#1},
  left=6pt,right=6pt,top=5pt,bottom=5pt}

%---------------------------------------------------------------------
% Code listings
%---------------------------------------------------------------------
\usepackage{upquote}      % straight ' and ` in listings, so copied code runs
\usepackage{attachfile2}  % embed scripts in the PDF, see \textattachfile
\lstdefinestyle{wkcode}{
  upquote=true,
  basicstyle=\ttfamily\footnotesize\color{weknowgraphite},
  backgroundcolor=\color{weknowgreyL},
  frame=single, framerule=0pt, rulecolor=\color{weknowgreyL},
  breaklines=true, breakatwhitespace=false,
  columns=fullflexible, keepspaces=true,
  aboveskip=8pt, belowskip=8pt, xleftmargin=6pt, framexleftmargin=6pt}
\lstset{style=wkcode}

\newcolumntype{L}[1]{>{\raggedright\arraybackslash}p{#1}}
\setlist{nosep,itemsep=4pt}

%---------------------------------------------------------------------
% Page-breaking behaviour
%   - never strand a single line at the top or bottom of a page
%   - \wkneed{N\baselineskip} makes a page break attractive when less
%     than N lines remain, so a heading cannot land at the page foot
%     with nothing under it (needspace's trick, inlined to avoid a
%     dependency)
%---------------------------------------------------------------------
\widowpenalty=10000
\clubpenalty=10000
\displaywidowpenalty=10000
\makeatletter
\newcommand{\wkneed}[1]{%
  \begingroup
    \setlength{\dimen@}{#1}%
    \vskip\z@\@plus\dimen@
    \penalty-100
    \vskip\z@\@plus-\dimen@
  \endgroup}
\makeatother

%---------------------------------------------------------------------
% Shorthands
%---------------------------------------------------------------------
\newcommand{\wpath}[1]{{\ttfamily\small #1}}
\newcommand{\tick}{{\color{weknowsea}$\checkmark$}}
\newcommand{\emptybox}{{\color{weknowgreyM}$\square$}}

% Role badges -- white on solid Azure, in two sizes.
%   \role     for the section heading (set \LARGE, so the badge sits at
%             \normalsize and reads as a tag, not a second title)
%   \roletoc  for the table of contents (set \normalsize, so the badge
%             drops to \scriptsize)
% Do not call these directly in a heading -- use \wksection or
% \wkrolesection below, which pick the right size per context.
% Robust: these are written to the .toc file, so they must not expand there.
% #1 size, #2 padding, #3 baseline raise, #4 text
\DeclareRobustCommand{\wkbadge}[4]{%
  {\setlength{\fboxsep}{#2}%
   \raisebox{#3}{%
     \colorbox{weknowazure}{\textcolor{white}{#1\bfseries #4}}}}}
\DeclareRobustCommand{\role}[1]{\texorpdfstring{%
  \hspace{0.45em}\wkbadge{\normalsize}{3.5pt}{-0.35ex}{#1}}{[#1]}}
\DeclareRobustCommand{\roletoc}[1]{\texorpdfstring{%
  \hspace{0.35em}\wkbadge{\scriptsize}{2pt}{-0.25ex}{#1}}{[#1]}}

%---------------------------------------------------------------------
% Section wrappers -- these set the running head and, where a role is
% given, place the badge at heading size in the body and at the smaller
% size in the table of contents.
%   \wksection{Final checklist}
%   \wkrolesection{Part 2 --- Install and lock the kiosk}{IT}
%---------------------------------------------------------------------
\newcommand{\wksection}[1]{%
  \renewcommand{\wksecmark}{#1}%
  \wkneed{6\baselineskip}%
  \section{#1}}
\newcommand{\wkrolesection}[2]{%
  \renewcommand{\wksecmark}{#1}%
  \wkneed{6\baselineskip}%
  \section[#1 \roletoc{#2}]{#1 \role{#2}}}

%---------------------------------------------------------------------
% Cover page -- \makewkgcover{title}{subtitle}{audience}
%---------------------------------------------------------------------
\newcommand{\makewkgcover}[3]{%
\begin{titlepage}
\thispagestyle{empty}
\begin{tikzpicture}[remember picture,overlay]
  % --- 'Wink' echo: symmetric Azure-tint rounded squares framing the top ---
  \fill[weknowazure20,rounded corners=8mm]
       ([xshift=-46mm,yshift=-46mm]current page.north east) rectangle
       ([xshift=14mm,yshift=-8mm]current page.north east);
  \fill[weknowazure40,rounded corners=5mm]
       ([xshift=-30mm,yshift=-40mm]current page.north east) rectangle
       ([xshift=8mm,yshift=-13mm]current page.north east);
  \fill[weknowazure20,rounded corners=8mm]
       ([xshift=-14mm,yshift=-46mm]current page.north west) rectangle
       ([xshift=46mm,yshift=-8mm]current page.north west);
  \fill[weknowazure40,rounded corners=5mm]
       ([xshift=-8mm,yshift=-40mm]current page.north west) rectangle
       ([xshift=30mm,yshift=-13mm]current page.north west);

  % --- Large centred WeKnow logo (stacked variant) = the hero ---
  \node[anchor=north] at ([yshift=-50mm]current page.north)
       {\includegraphics[width=52mm]{wkg-logo-stacked}};

  % --- Centred title / Azure hero rule / subtitle ---
  \node[anchor=north,align=center,text width=185mm]
       at ([yshift=-126mm]current page.north)
       {{\color{weknowgraphite}\fontsize{26}{31}\selectfont\bfseries #1}\\[11pt]
        {\color{weknowazure}\rule{42mm}{2pt}}\\[11pt]
        {\color{weknowsea}\fontsize{15}{19}\selectfont #2}};

  % --- Metadata panel ---
  \node[anchor=north, fill=weknowsky, rounded corners=3mm,
        inner xsep=12mm, inner ysep=8mm]
       at ([yshift=-178mm]current page.north)
       {\renewcommand{\arraystretch}{1.5}%
        \begin{tabular}{@{}l@{\hspace{12mm}}l@{}}
          {\color{weknowsea}\bfseries Version} & {\color{weknowgraphite}1.0}\\
          {\color{weknowsea}\bfseries Date} & {\color{weknowgraphite}August 2026}\\
          {\color{weknowsea}\bfseries Classification} & {\color{weknowgraphite}Internal}\\
          {\color{weknowsea}\bfseries Prepared by} & {\color{weknowgraphite}Vedant Kalbag}\\
          {\color{weknowsea}\bfseries For} & {\color{weknowgraphite}#3}\\
        \end{tabular}};

  % --- Confidentiality notice ---
  \node[anchor=south,align=center,text width=150mm]
       at ([yshift=20mm]current page.south)
       {{\color{weknowgreyM}\rule{55mm}{0.4pt}}\\[8pt]
        {\color{weknowgreyN}\footnotesize This document contains confidential and proprietary
         information belonging to We Know Group. It is intended solely for internal use and may
         not be reproduced or distributed without prior written consent.}};
\end{tikzpicture}
\end{titlepage}}

%      \renewcommand{\guidetitle}{Short title for the page footer}
%=====================================================================

%---------------------------------------------------------------------
% Fonts & encoding -- Inter (brand stand-in for Circular Pro)
%---------------------------------------------------------------------
\usepackage[T1]{fontenc}
\usepackage[utf8]{inputenc}
\usepackage[sfdefault]{inter}
\renewcommand{\familydefault}{\sfdefault}
\usepackage{microtype}

%---------------------------------------------------------------------
% Layout
%---------------------------------------------------------------------
\usepackage[a4paper,margin=20mm,top=28mm,bottom=22mm,headheight=20pt,headsep=14pt]{geometry}
\usepackage{graphicx}
\graphicspath{{images/}}
\usepackage[table]{xcolor}
\usepackage{tikz}
\usetikzlibrary{arrows.meta,positioning,calc,shapes.misc,fit,backgrounds}
\usepackage{tcolorbox}
\tcbuselibrary{skins,breakable}
\usepackage{booktabs}
\usepackage{array}
\usepackage{tabularx}
\usepackage{enumitem}
\usepackage{fancyhdr}
\usepackage{titlesec}
\usepackage{listings}
\usepackage{float}
\usepackage{lastpage}
\usepackage{amssymb}
\usepackage{caption}
\usepackage[hidelinks,colorlinks=true,urlcolor=weknowsea,linkcolor=weknownight]{hyperref}

%---------------------------------------------------------------------
% We Know Group palette (Brand Guidelines V2 2023, Section 3)
%---------------------------------------------------------------------
\definecolor{weknowazure}{HTML}{00A6D3}     % Hero colour
\definecolor{weknowgraphite}{HTML}{121212}
\definecolor{weknowsky}{HTML}{E5F1F9}
\definecolor{weknowsea}{HTML}{008BB2}
\definecolor{weknownight}{HTML}{0C2752}
\definecolor{weknowgreyN}{HTML}{575A5C}
\definecolor{weknowgreyM}{HTML}{ADADAD}
\definecolor{weknowgreyL}{HTML}{F4F4F4}
\definecolor{weknowazure80}{HTML}{33B8DC}    % Azure tints (brand: tint Azure only)
\definecolor{weknowazure60}{HTML}{66CAE5}
\definecolor{weknowazure40}{HTML}{99DBED}
\definecolor{weknowazure20}{HTML}{CCEDF6}

%---------------------------------------------------------------------
% Headers / footers -- \guidetitle is set by each document
%---------------------------------------------------------------------
\providecommand{\guidetitle}{Kato Kiosk Guide}
\pagestyle{fancy}
\fancyhf{}
% \sectionmark uses \markright, so the running head must read \rightmark.
\fancyhead[L]{\footnotesize\color{weknowgreyN}\nouppercase{\rightmark}}
\fancyhead[R]{\raisebox{-0.35\height}{\includegraphics[height=15pt]{wkg-logo-horizontal}}}
\renewcommand{\headrulewidth}{0.8pt}
\renewcommand{\headrule}{\hbox to\headwidth{\color{weknowazure}\leaders\hrule height \headrulewidth\hfill}}
\fancyfoot[L]{\scriptsize\color{weknowgreyN}Internal --- We Know Group}
\fancyfoot[C]{\scriptsize\color{weknowgreyN}\guidetitle}
\fancyfoot[R]{\scriptsize\color{weknowgreyN}Page \thepage\ of \pageref{LastPage}}
\renewcommand{\footrulewidth}{0.4pt}
\renewcommand{\footrule}{\hbox to\headwidth{\color{weknowgreyM}\leaders\hrule height \footrulewidth\hfill}}

%---------------------------------------------------------------------
% Section styling
%---------------------------------------------------------------------
\titleformat{\section}
  {\color{weknowgraphite}\LARGE\bfseries}
  {\color{weknowazure}\thesection}{0.7em}{}
  [{\vspace{3pt}\color{weknowazure}\titlerule[1.4pt]}]
\titleformat{\subsection}
  {\color{weknowsea}\Large\bfseries}
  {\color{weknowazure}\thesubsection}{0.6em}{}
\titlespacing*{\section}{0pt}{20pt}{12pt}
\titlespacing*{\subsection}{0pt}{16pt}{6pt}
% The running head takes its text from \wksecmark, not from the section
% title, so role badges never leak into the header. \wksection and
% \wkrolesection (defined below) set it. Subsections must not overwrite
% the section in the head.
\newcommand{\wksecmark}{}
\renewcommand{\sectionmark}[1]{\markright{\thesection\quad\wksecmark}}
\renewcommand{\subsectionmark}[1]{}

%---------------------------------------------------------------------
% Captions
%---------------------------------------------------------------------
\captionsetup{font=small,labelfont={bf,color=weknownight},
              textfont={color=weknowgreyN},skip=5pt}

%---------------------------------------------------------------------
% Screenshot helper -- \screenshot[width fraction]{file}{caption}
%---------------------------------------------------------------------
\newcommand{\screenshot}[3][0.72]{%
  \begin{figure}[H]
    \centering
    \setlength{\fboxsep}{0pt}\setlength{\fboxrule}{0.6pt}%
    {\color{weknowgreyM}\fbox{\includegraphics[width=#1\linewidth]{#2}}}%
    \caption{#3}
  \end{figure}}

%---------------------------------------------------------------------
% Callout boxes
%---------------------------------------------------------------------
\newtcolorbox{wkinfo}[1][]{%
  enhanced,breakable,colback=weknowsky,colframe=weknowazure,
  boxrule=0pt,leftrule=2.5pt,arc=1.5pt,
  fonttitle=\bfseries,coltitle=weknownight,title={#1},
  left=6pt,right=6pt,top=5pt,bottom=5pt}
\newtcolorbox{wkwarn}[1][]{%
  enhanced,breakable,colback=weknowgreyL,colframe=weknownight,
  boxrule=0pt,leftrule=2.5pt,arc=1.5pt,
  fonttitle=\bfseries,coltitle=white,title={#1},
  left=6pt,right=6pt,top=5pt,bottom=5pt}

%---------------------------------------------------------------------
% Code listings
%---------------------------------------------------------------------
\usepackage{upquote}      % straight ' and ` in listings, so copied code runs
\usepackage{attachfile2}  % embed scripts in the PDF, see \textattachfile
\lstdefinestyle{wkcode}{
  upquote=true,
  basicstyle=\ttfamily\footnotesize\color{weknowgraphite},
  backgroundcolor=\color{weknowgreyL},
  frame=single, framerule=0pt, rulecolor=\color{weknowgreyL},
  breaklines=true, breakatwhitespace=false,
  columns=fullflexible, keepspaces=true,
  aboveskip=8pt, belowskip=8pt, xleftmargin=6pt, framexleftmargin=6pt}
\lstset{style=wkcode}

\newcolumntype{L}[1]{>{\raggedright\arraybackslash}p{#1}}
\setlist{nosep,itemsep=4pt}

%---------------------------------------------------------------------
% Page-breaking behaviour
%   - never strand a single line at the top or bottom of a page
%   - \wkneed{N\baselineskip} makes a page break attractive when less
%     than N lines remain, so a heading cannot land at the page foot
%     with nothing under it (needspace's trick, inlined to avoid a
%     dependency)
%---------------------------------------------------------------------
\widowpenalty=10000
\clubpenalty=10000
\displaywidowpenalty=10000
\makeatletter
\newcommand{\wkneed}[1]{%
  \begingroup
    \setlength{\dimen@}{#1}%
    \vskip\z@\@plus\dimen@
    \penalty-100
    \vskip\z@\@plus-\dimen@
  \endgroup}
\makeatother

%---------------------------------------------------------------------
% Shorthands
%---------------------------------------------------------------------
\newcommand{\wpath}[1]{{\ttfamily\small #1}}
\newcommand{\tick}{{\color{weknowsea}$\checkmark$}}
\newcommand{\emptybox}{{\color{weknowgreyM}$\square$}}

% Role badges -- white on solid Azure, in two sizes.
%   \role     for the section heading (set \LARGE, so the badge sits at
%             \normalsize and reads as a tag, not a second title)
%   \roletoc  for the table of contents (set \normalsize, so the badge
%             drops to \scriptsize)
% Do not call these directly in a heading -- use \wksection or
% \wkrolesection below, which pick the right size per context.
% Robust: these are written to the .toc file, so they must not expand there.
% #1 size, #2 padding, #3 baseline raise, #4 text
\DeclareRobustCommand{\wkbadge}[4]{%
  {\setlength{\fboxsep}{#2}%
   \raisebox{#3}{%
     \colorbox{weknowazure}{\textcolor{white}{#1\bfseries #4}}}}}
\DeclareRobustCommand{\role}[1]{\texorpdfstring{%
  \hspace{0.45em}\wkbadge{\normalsize}{3.5pt}{-0.35ex}{#1}}{[#1]}}
\DeclareRobustCommand{\roletoc}[1]{\texorpdfstring{%
  \hspace{0.35em}\wkbadge{\scriptsize}{2pt}{-0.25ex}{#1}}{[#1]}}

%---------------------------------------------------------------------
% Section wrappers -- these set the running head and, where a role is
% given, place the badge at heading size in the body and at the smaller
% size in the table of contents.
%   \wksection{Final checklist}
%   \wkrolesection{Part 2 --- Install and lock the kiosk}{IT}
%---------------------------------------------------------------------
\newcommand{\wksection}[1]{%
  \renewcommand{\wksecmark}{#1}%
  \wkneed{6\baselineskip}%
  \section{#1}}
\newcommand{\wkrolesection}[2]{%
  \renewcommand{\wksecmark}{#1}%
  \wkneed{6\baselineskip}%
  \section[#1 \roletoc{#2}]{#1 \role{#2}}}

%---------------------------------------------------------------------
% Cover page -- \makewkgcover{title}{subtitle}{audience}
%---------------------------------------------------------------------
\newcommand{\makewkgcover}[3]{%
\begin{titlepage}
\thispagestyle{empty}
\begin{tikzpicture}[remember picture,overlay]
  % --- 'Wink' echo: symmetric Azure-tint rounded squares framing the top ---
  \fill[weknowazure20,rounded corners=8mm]
       ([xshift=-46mm,yshift=-46mm]current page.north east) rectangle
       ([xshift=14mm,yshift=-8mm]current page.north east);
  \fill[weknowazure40,rounded corners=5mm]
       ([xshift=-30mm,yshift=-40mm]current page.north east) rectangle
       ([xshift=8mm,yshift=-13mm]current page.north east);
  \fill[weknowazure20,rounded corners=8mm]
       ([xshift=-14mm,yshift=-46mm]current page.north west) rectangle
       ([xshift=46mm,yshift=-8mm]current page.north west);
  \fill[weknowazure40,rounded corners=5mm]
       ([xshift=-8mm,yshift=-40mm]current page.north west) rectangle
       ([xshift=30mm,yshift=-13mm]current page.north west);

  % --- Large centred WeKnow logo (stacked variant) = the hero ---
  \node[anchor=north] at ([yshift=-50mm]current page.north)
       {\includegraphics[width=52mm]{wkg-logo-stacked}};

  % --- Centred title / Azure hero rule / subtitle ---
  \node[anchor=north,align=center,text width=185mm]
       at ([yshift=-126mm]current page.north)
       {{\color{weknowgraphite}\fontsize{26}{31}\selectfont\bfseries #1}\\[11pt]
        {\color{weknowazure}\rule{42mm}{2pt}}\\[11pt]
        {\color{weknowsea}\fontsize{15}{19}\selectfont #2}};

  % --- Metadata panel ---
  \node[anchor=north, fill=weknowsky, rounded corners=3mm,
        inner xsep=12mm, inner ysep=8mm]
       at ([yshift=-178mm]current page.north)
       {\renewcommand{\arraystretch}{1.5}%
        \begin{tabular}{@{}l@{\hspace{12mm}}l@{}}
          {\color{weknowsea}\bfseries Version} & {\color{weknowgraphite}1.0}\\
          {\color{weknowsea}\bfseries Date} & {\color{weknowgraphite}August 2026}\\
          {\color{weknowsea}\bfseries Classification} & {\color{weknowgraphite}Internal}\\
          {\color{weknowsea}\bfseries Prepared by} & {\color{weknowgraphite}Vedant Kalbag}\\
          {\color{weknowsea}\bfseries For} & {\color{weknowgraphite}#3}\\
        \end{tabular}};

  % --- Confidentiality notice ---
  \node[anchor=south,align=center,text width=150mm]
       at ([yshift=20mm]current page.south)
       {{\color{weknowgreyM}\rule{55mm}{0.4pt}}\\[8pt]
        {\color{weknowgreyN}\footnotesize This document contains confidential and proprietary
         information belonging to We Know Group. It is intended solely for internal use and may
         not be reproduced or distributed without prior written consent.}};
\end{tikzpicture}
\end{titlepage}}

%      \renewcommand{\guidetitle}{Short title for the page footer}
%=====================================================================

%---------------------------------------------------------------------
% Fonts & encoding -- Inter (brand stand-in for Circular Pro)
%---------------------------------------------------------------------
\usepackage[T1]{fontenc}
\usepackage[utf8]{inputenc}
\usepackage[sfdefault]{inter}
\renewcommand{\familydefault}{\sfdefault}
\usepackage{microtype}

%---------------------------------------------------------------------
% Layout
%---------------------------------------------------------------------
\usepackage[a4paper,margin=20mm,top=28mm,bottom=22mm,headheight=20pt,headsep=14pt]{geometry}
\usepackage{graphicx}
\graphicspath{{images/}}
\usepackage[table]{xcolor}
\usepackage{tikz}
\usetikzlibrary{arrows.meta,positioning,calc,shapes.misc,fit,backgrounds}
\usepackage{tcolorbox}
\tcbuselibrary{skins,breakable}
\usepackage{booktabs}
\usepackage{array}
\usepackage{tabularx}
\usepackage{enumitem}
\usepackage{fancyhdr}
\usepackage{titlesec}
\usepackage{listings}
\usepackage{float}
\usepackage{lastpage}
\usepackage{amssymb}
\usepackage{caption}
\usepackage[hidelinks,colorlinks=true,urlcolor=weknowsea,linkcolor=weknownight]{hyperref}

%---------------------------------------------------------------------
% We Know Group palette (Brand Guidelines V2 2023, Section 3)
%---------------------------------------------------------------------
\definecolor{weknowazure}{HTML}{00A6D3}     % Hero colour
\definecolor{weknowgraphite}{HTML}{121212}
\definecolor{weknowsky}{HTML}{E5F1F9}
\definecolor{weknowsea}{HTML}{008BB2}
\definecolor{weknownight}{HTML}{0C2752}
\definecolor{weknowgreyN}{HTML}{575A5C}
\definecolor{weknowgreyM}{HTML}{ADADAD}
\definecolor{weknowgreyL}{HTML}{F4F4F4}
\definecolor{weknowazure80}{HTML}{33B8DC}    % Azure tints (brand: tint Azure only)
\definecolor{weknowazure60}{HTML}{66CAE5}
\definecolor{weknowazure40}{HTML}{99DBED}
\definecolor{weknowazure20}{HTML}{CCEDF6}

%---------------------------------------------------------------------
% Headers / footers -- \guidetitle is set by each document
%---------------------------------------------------------------------
\providecommand{\guidetitle}{Kato Kiosk Guide}
\pagestyle{fancy}
\fancyhf{}
% \sectionmark uses \markright, so the running head must read \rightmark.
\fancyhead[L]{\footnotesize\color{weknowgreyN}\nouppercase{\rightmark}}
\fancyhead[R]{\raisebox{-0.35\height}{\includegraphics[height=15pt]{wkg-logo-horizontal}}}
\renewcommand{\headrulewidth}{0.8pt}
\renewcommand{\headrule}{\hbox to\headwidth{\color{weknowazure}\leaders\hrule height \headrulewidth\hfill}}
\fancyfoot[L]{\scriptsize\color{weknowgreyN}Internal --- We Know Group}
\fancyfoot[C]{\scriptsize\color{weknowgreyN}\guidetitle}
\fancyfoot[R]{\scriptsize\color{weknowgreyN}Page \thepage\ of \pageref{LastPage}}
\renewcommand{\footrulewidth}{0.4pt}
\renewcommand{\footrule}{\hbox to\headwidth{\color{weknowgreyM}\leaders\hrule height \footrulewidth\hfill}}

%---------------------------------------------------------------------
% Section styling
%---------------------------------------------------------------------
\titleformat{\section}
  {\color{weknowgraphite}\LARGE\bfseries}
  {\color{weknowazure}\thesection}{0.7em}{}
  [{\vspace{3pt}\color{weknowazure}\titlerule[1.4pt]}]
\titleformat{\subsection}
  {\color{weknowsea}\Large\bfseries}
  {\color{weknowazure}\thesubsection}{0.6em}{}
\titlespacing*{\section}{0pt}{20pt}{12pt}
\titlespacing*{\subsection}{0pt}{16pt}{6pt}
% The running head takes its text from \wksecmark, not from the section
% title, so role badges never leak into the header. \wksection and
% \wkrolesection (defined below) set it. Subsections must not overwrite
% the section in the head.
\newcommand{\wksecmark}{}
\renewcommand{\sectionmark}[1]{\markright{\thesection\quad\wksecmark}}
\renewcommand{\subsectionmark}[1]{}

%---------------------------------------------------------------------
% Captions
%---------------------------------------------------------------------
\captionsetup{font=small,labelfont={bf,color=weknownight},
              textfont={color=weknowgreyN},skip=5pt}

%---------------------------------------------------------------------
% Screenshot helper -- \screenshot[width fraction]{file}{caption}
%---------------------------------------------------------------------
\newcommand{\screenshot}[3][0.72]{%
  \begin{figure}[H]
    \centering
    \setlength{\fboxsep}{0pt}\setlength{\fboxrule}{0.6pt}%
    {\color{weknowgreyM}\fbox{\includegraphics[width=#1\linewidth]{#2}}}%
    \caption{#3}
  \end{figure}}

%---------------------------------------------------------------------
% Callout boxes
%---------------------------------------------------------------------
\newtcolorbox{wkinfo}[1][]{%
  enhanced,breakable,colback=weknowsky,colframe=weknowazure,
  boxrule=0pt,leftrule=2.5pt,arc=1.5pt,
  fonttitle=\bfseries,coltitle=weknownight,title={#1},
  left=6pt,right=6pt,top=5pt,bottom=5pt}
\newtcolorbox{wkwarn}[1][]{%
  enhanced,breakable,colback=weknowgreyL,colframe=weknownight,
  boxrule=0pt,leftrule=2.5pt,arc=1.5pt,
  fonttitle=\bfseries,coltitle=white,title={#1},
  left=6pt,right=6pt,top=5pt,bottom=5pt}

%---------------------------------------------------------------------
% Code listings
%---------------------------------------------------------------------
\usepackage{upquote}      % straight ' and ` in listings, so copied code runs
\usepackage{attachfile2}  % embed scripts in the PDF, see \textattachfile
\lstdefinestyle{wkcode}{
  upquote=true,
  basicstyle=\ttfamily\footnotesize\color{weknowgraphite},
  backgroundcolor=\color{weknowgreyL},
  frame=single, framerule=0pt, rulecolor=\color{weknowgreyL},
  breaklines=true, breakatwhitespace=false,
  columns=fullflexible, keepspaces=true,
  aboveskip=8pt, belowskip=8pt, xleftmargin=6pt, framexleftmargin=6pt}
\lstset{style=wkcode}

\newcolumntype{L}[1]{>{\raggedright\arraybackslash}p{#1}}
\setlist{nosep,itemsep=4pt}

%---------------------------------------------------------------------
% Page-breaking behaviour
%   - never strand a single line at the top or bottom of a page
%   - \wkneed{N\baselineskip} makes a page break attractive when less
%     than N lines remain, so a heading cannot land at the page foot
%     with nothing under it (needspace's trick, inlined to avoid a
%     dependency)
%---------------------------------------------------------------------
\widowpenalty=10000
\clubpenalty=10000
\displaywidowpenalty=10000
\makeatletter
\newcommand{\wkneed}[1]{%
  \begingroup
    \setlength{\dimen@}{#1}%
    \vskip\z@\@plus\dimen@
    \penalty-100
    \vskip\z@\@plus-\dimen@
  \endgroup}
\makeatother

%---------------------------------------------------------------------
% Shorthands
%---------------------------------------------------------------------
\newcommand{\wpath}[1]{{\ttfamily\small #1}}
\newcommand{\tick}{{\color{weknowsea}$\checkmark$}}
\newcommand{\emptybox}{{\color{weknowgreyM}$\square$}}

% Role badges -- white on solid Azure, in two sizes.
%   \role     for the section heading (set \LARGE, so the badge sits at
%             \normalsize and reads as a tag, not a second title)
%   \roletoc  for the table of contents (set \normalsize, so the badge
%             drops to \scriptsize)
% Do not call these directly in a heading -- use \wksection or
% \wkrolesection below, which pick the right size per context.
% Robust: these are written to the .toc file, so they must not expand there.
% #1 size, #2 padding, #3 baseline raise, #4 text
\DeclareRobustCommand{\wkbadge}[4]{%
  {\setlength{\fboxsep}{#2}%
   \raisebox{#3}{%
     \colorbox{weknowazure}{\textcolor{white}{#1\bfseries #4}}}}}
\DeclareRobustCommand{\role}[1]{\texorpdfstring{%
  \hspace{0.45em}\wkbadge{\normalsize}{3.5pt}{-0.35ex}{#1}}{[#1]}}
\DeclareRobustCommand{\roletoc}[1]{\texorpdfstring{%
  \hspace{0.35em}\wkbadge{\scriptsize}{2pt}{-0.25ex}{#1}}{[#1]}}

%---------------------------------------------------------------------
% Section wrappers -- these set the running head and, where a role is
% given, place the badge at heading size in the body and at the smaller
% size in the table of contents.
%   \wksection{Final checklist}
%   \wkrolesection{Part 2 --- Install and lock the kiosk}{IT}
%---------------------------------------------------------------------
\newcommand{\wksection}[1]{%
  \renewcommand{\wksecmark}{#1}%
  \wkneed{6\baselineskip}%
  \section{#1}}
\newcommand{\wkrolesection}[2]{%
  \renewcommand{\wksecmark}{#1}%
  \wkneed{6\baselineskip}%
  \section[#1 \roletoc{#2}]{#1 \role{#2}}}

%---------------------------------------------------------------------
% Cover page -- \makewkgcover{title}{subtitle}{audience}
%---------------------------------------------------------------------
\newcommand{\makewkgcover}[3]{%
\begin{titlepage}
\thispagestyle{empty}
\begin{tikzpicture}[remember picture,overlay]
  % --- 'Wink' echo: symmetric Azure-tint rounded squares framing the top ---
  \fill[weknowazure20,rounded corners=8mm]
       ([xshift=-46mm,yshift=-46mm]current page.north east) rectangle
       ([xshift=14mm,yshift=-8mm]current page.north east);
  \fill[weknowazure40,rounded corners=5mm]
       ([xshift=-30mm,yshift=-40mm]current page.north east) rectangle
       ([xshift=8mm,yshift=-13mm]current page.north east);
  \fill[weknowazure20,rounded corners=8mm]
       ([xshift=-14mm,yshift=-46mm]current page.north west) rectangle
       ([xshift=46mm,yshift=-8mm]current page.north west);
  \fill[weknowazure40,rounded corners=5mm]
       ([xshift=-8mm,yshift=-40mm]current page.north west) rectangle
       ([xshift=30mm,yshift=-13mm]current page.north west);

  % --- Large centred WeKnow logo (stacked variant) = the hero ---
  \node[anchor=north] at ([yshift=-50mm]current page.north)
       {\includegraphics[width=52mm]{wkg-logo-stacked}};

  % --- Centred title / Azure hero rule / subtitle ---
  \node[anchor=north,align=center,text width=185mm]
       at ([yshift=-126mm]current page.north)
       {{\color{weknowgraphite}\fontsize{26}{31}\selectfont\bfseries #1}\\[11pt]
        {\color{weknowazure}\rule{42mm}{2pt}}\\[11pt]
        {\color{weknowsea}\fontsize{15}{19}\selectfont #2}};

  % --- Metadata panel ---
  \node[anchor=north, fill=weknowsky, rounded corners=3mm,
        inner xsep=12mm, inner ysep=8mm]
       at ([yshift=-178mm]current page.north)
       {\renewcommand{\arraystretch}{1.5}%
        \begin{tabular}{@{}l@{\hspace{12mm}}l@{}}
          {\color{weknowsea}\bfseries Version} & {\color{weknowgraphite}1.0}\\
          {\color{weknowsea}\bfseries Date} & {\color{weknowgraphite}August 2026}\\
          {\color{weknowsea}\bfseries Classification} & {\color{weknowgraphite}Internal}\\
          {\color{weknowsea}\bfseries Prepared by} & {\color{weknowgraphite}Vedant Kalbag}\\
          {\color{weknowsea}\bfseries For} & {\color{weknowgraphite}#3}\\
        \end{tabular}};

  % --- Confidentiality notice ---
  \node[anchor=south,align=center,text width=150mm]
       at ([yshift=20mm]current page.south)
       {{\color{weknowgreyM}\rule{55mm}{0.4pt}}\\[8pt]
        {\color{weknowgreyN}\footnotesize This document contains confidential and proprietary
         information belonging to We Know Group. It is intended solely for internal use and may
         not be reproduced or distributed without prior written consent.}};
\end{tikzpicture}
\end{titlepage}}

\renewcommand{\guidetitle}{Kato Kiosk Configuration Guide --- AM \& IT}

\begin{document}

\makewkgcover{Kato:\\ Kiosk Configuration Guide}
             {Configuring a kiosk for installation: \\Command Centre, App Install \& Lockdown}
             {Account Managers \& IT}

%=====================================================================
% HOW TO USE
%=====================================================================
\section*{How to use this guide}
\markright{How to use this guide}

This guide covers \textbf{configuring a kiosk for a specific
installation}. It has two parts, split by who does the work:

\begin{center}
\renewcommand{\arraystretch}{1.35}\small
\begin{tabularx}{\textwidth}{L{20mm} L{34mm} X}
\rowcolor{weknownight}\color{white}\textbf{Part} &
\color{white}\textbf{Who} & \color{white}\textbf{What}\\
\textbf{Part 1} & Account Manager &
Create the location, kiosk and installation record in Command Centre,
clear the prep checks, and hand the Kiosk ID and API key to IT.\\
\rowcolor{weknowgreyL}
\textbf{Part 2} & IT &
Install the Kato app on the machine using those credentials, and lock the
machine into kiosk mode.\\
\end{tabularx}
\end{center}

\begin{figure}[H]
\centering
\begin{tikzpicture}[
  box/.style={draw=weknowazure,fill=white,rounded corners=3pt,
              text width=26mm,minimum height=12mm,align=center,
              inner sep=3pt,font=\scriptsize},
  lane/.style={font=\small\bfseries,text=weknowsea,align=center,
               text width=20mm},
  arr/.style={-{Stealth[length=2.4mm]},thick,weknowazure},
  handoff/.style={-{Stealth[length=2.4mm]},thick,weknowgreyN,dashed}]

  % --- Account Manager lane ---
  \node[box]  (t1) at (0,0)   {Create location\\and kiosk};
  \node[box]  (t2) at (3.5,0) {Configure and\\clear prep checks};
  \node[box]  (t3) at (7.0,0) {Generate\\credentials};
  \node[box]  (t4) at (10.5,0) {Create installation\\record};
  \node[lane] (l1) at (-2.8,0) {Part 1\\Account Manager};

  % --- IT lane ---
  \node[box]  (b1) at (7.0,-2.7) {Install the\\Kato app};
  \node[box]  (b2) at (10.5,-2.7) {Lock into\\kiosk mode};
  \node[lane] (l2) at (-2.8,-2.7) {Part 2\\IT};

  % --- Lane bands ---
  \begin{scope}[on background layer]
    \node[fill=weknowsky,rounded corners=4pt,inner sep=3.5mm,
          fit=(l1)(t1)(t4)] {};
    \node[fill=weknowgreyL,rounded corners=4pt,inner sep=3.5mm,
          fit=(l2)(b1)(b2)] {};
  \end{scope}

  \draw[arr] (t1) -- (t2);
  \draw[arr] (t2) -- (t3);
  \draw[arr] (t3) -- (t4);
  \draw[arr] (b1) -- (b2);

  \draw[handoff] (t3.south) -- (b1.north);
  \node[font=\scriptsize,text=weknowgreyN,anchor=west]
       at ([xshift=2.5mm]$(t3.south)!0.5!(b1.north)$) {Kiosk ID + API key};
\end{tikzpicture}
\caption{The Account Manager configures the kiosk and hands two credentials to IT, who install the app and lock the machine down.}
\end{figure}

\begin{wkwarn}[The machine must already be built]
This guide assumes FreedomPay FCC and the PAL packages are installed and
verified on the machine. If that has not been done, work through the
\emph{Kato Kiosk Setup Guide} first --- the prep checks in Part~1 cannot
honestly be cleared until it is.
\end{wkwarn}

\begin{wkinfo}[Golden rules]
\begin{itemize}
  \item The \textbf{API key is shown once}. Copy it the moment you
        generate it, or you will have to generate a new one.
  \item The \textbf{Kiosk ID} and the \textbf{Workstation ID} are
        different things --- see \S\ref{sec:ids}.
  \item Do not clear a prep check you have not actually verified. It is
        the only thing standing between a half-built kiosk and a live
        one.
\end{itemize}
\end{wkinfo}

\clearpage
\tableofcontents
\clearpage

%=====================================================================
% 1  BEFORE YOU START
%=====================================================================
\wksection{Before you start}

\subsection{Minimum information needed for a location}

The Account Manager must have all of the following for the location
before an installation can be created. If any are missing, chase them
first.

\begin{center}
\renewcommand{\arraystretch}{1.3}\small
\begin{tabularx}{\textwidth}{X X}
\rowcolor{weknowgreyL}
\emptybox\ Hotel name & \emptybox\ Currency\\
\emptybox\ Which products to enable & \emptybox\ Timezone\\
\rowcolor{weknowgreyL}
\emptybox\ Products for the home screen & \emptybox\ Relevant airports\\
\emptybox\ Default language (kiosk + PDQ) & \emptybox\ Default Uber suggestions\\
\rowcolor{weknowgreyL}
\emptybox\ VAT / Tax ID & \emptybox\ Tax entity name\\
\emptybox\ Banking details (in full) & \\
\end{tabularx}
\end{center}

This is the \emph{minimum}. Many other parameters can be configured at
each level of the hierarchy --- see \S\ref{sec:hierarchy} for where each
one belongs.

\begin{wkwarn}[Legacy fields --- not required]
\textbf{RPS customer code} and \textbf{outlet code} are legacy fields.
Do not chase the site for them.
\end{wkwarn}

\subsection{Two identifiers, two different things}\label{sec:ids}

\begin{center}
\renewcommand{\arraystretch}{1.35}\small
\begin{tabularx}{\textwidth}{L{34mm} X}
\rowcolor{weknownight}\color{white}\textbf{Identifier} &
\color{white}\textbf{What it is}\\
\textbf{Kiosk ID} &
A UUID generated by Command Centre when the kiosk record is created. The
app uses it to identify itself to the server.\\
\rowcolor{weknowgreyL}
\textbf{Workstation ID} &
The physical machine label, e.g.\ \wpath{WKG-POS-617}. Also entered as the
kiosk's serial number in Command Centre.\\
\end{tabularx}
\end{center}

Both go into the install command in Part~2 --- do not swap them.

\subsection{Important paths}\label{sec:paths}

Every file and setting this guide touches on the kiosk machine.

\begin{center}
\renewcommand{\arraystretch}{1.3}\small
\begin{tabularx}{\textwidth}{L{34mm} X}
\rowcolor{weknownight}\color{white}\textbf{What} &
\color{white}\textbf{Where}\\
Identity file &
\wpath{C:\textbackslash Users\textbackslash SelfBooker\textbackslash AppData\textbackslash Roaming\textbackslash WKG Kiosk\textbackslash kiosk-identity.json}\\
\rowcolor{weknowgreyL}
Kato app &
\wpath{C:\textbackslash Users\textbackslash SelfBooker\textbackslash AppData\textbackslash Local\textbackslash Programs\textbackslash Kiosk\textbackslash Kiosk.exe}\\
Installer drop folder &
\wpath{C:\textbackslash Users\textbackslash SelfBooker\textbackslash Documents}\\
\rowcolor{weknowgreyL}
Shell Launcher utility &
\wpath{C:\textbackslash Kiosk\textbackslash SSM25\_files\textbackslash Shell-Launcher-Utilities\textbackslash SLUtility.exe}\\
FCC logs &
\wpath{C:\textbackslash ProgramData\textbackslash FreedomPay\textbackslash Freeway Commerce Connect\textbackslash log}\\
\rowcolor{weknowgreyL}
FCC registry key &
\wpath{HKEY\_LOCAL\_MACHINE\textbackslash SOFTWARE\textbackslash Freedompay\textbackslash Atlas\textbackslash PAL}\\
Latest app release &
\href{https://github.com/WeKnowTechnology/wkg-electron-app/releases}{\wpath{github.com/WeKnowTechnology/wkg-electron-app/releases}}\\
\end{tabularx}
\end{center}

%=====================================================================
% 2  CONFIGURATION HIERARCHY
%=====================================================================
\wkrolesection{How configuration resolves}{Account Manager}\label{sec:hierarchy}

Configuration is layered. An individual kiosk should carry the
\textbf{minimum configuration needed} --- only the things genuinely unique
to that kiosk. Everything else resolves down from the levels above it.

\begin{figure}[H]
\centering
\begin{tikzpicture}[
  lvl/.style={rounded corners=2.5pt,draw=weknowazure60,
              minimum height=8mm,align=center,inner xsep=4pt,
              font=\scriptsize,text=weknownight},
  node distance=1.6mm]

  \node[lvl,fill=weknowgreyL]    (g)  {Global};
  \node[lvl,fill=weknowazure20,right=of g]  (co) {Country};
  \node[lvl,fill=weknowazure20,right=of co] (re) {Region\\(State)};
  \node[lvl,fill=weknowazure40,right=of re] (ci) {City};
  \node[lvl,fill=weknowazure40,right=of ci] (lg) {Location\\Group};
  \node[lvl,fill=weknowazure60,right=of lg] (lo) {Location};
  \node[lvl,fill=weknowazure60,right=of lo] (kg) {Kiosk\\Group};
  \node[lvl,fill=weknowazure80,right=of kg] (ki) {Kiosk};

  \draw[-{Stealth[length=2.4mm]},thick,weknowazure]
       ([yshift=-5mm]g.south west) -- ([yshift=-5mm]ki.south east);
  \node[font=\scriptsize,text=weknowgreyN,anchor=north]
       at ([yshift=-6mm]$(g.south west)!0.5!(ki.south east)$)
       {widest scope \quad$\longrightarrow$\quad most specific --- and highest priority};
\end{tikzpicture}
\caption{The configuration hierarchy. Anything set further right overrides the same parameter set further left.}
\end{figure}

\subsection{The rule}

A parameter set at a \textbf{more granular} level overrides the same
parameter set at a \textbf{wider} level. So a currency set on a Country
applies to every location in it, until a specific Location --- or a single
Kiosk --- sets its own.

\begin{wkinfo}[Set each parameter at the widest level where it holds true]
If every kiosk in a country shares a currency, set it once on the Country;
do not repeat it on each kiosk. Repeating a value further down does no
harm on day one, but it is the thing that makes later changes miss ---
you update the Country and the kiosk keeps its own stale override.
\end{wkinfo}

\subsection{When a kiosk is in more than one group}\label{sec:groupsorder}

A kiosk can belong to \textbf{several kiosk groups at once}. Assign them
on the kiosk's \textbf{Edit} tab, under \textbf{Group assignments}.

Assigned groups carry a \textbf{display order}, and that order decides
conflicts: where two of a kiosk's groups set the same parameter, the group
\textbf{higher in the list wins} and overrides the ones below it.

\begin{wkinfo}[Order the groups deliberately]
Put the most specific group at the top. A kiosk that is both
\emph{eSIM-only} and part of a wider \emph{Heathrow} group should carry
the narrower group above the broader one, so the narrow group's settings
are the ones that survive.
\end{wkinfo}

\subsection{Working out where a parameter belongs}

Ask how widely the value is true, and set it at that level:

\begin{center}
\renewcommand{\arraystretch}{1.35}\small
\begin{tabularx}{\textwidth}{L{34mm} X}
\rowcolor{weknownight}\color{white}\textbf{Level} &
\color{white}\textbf{Typically carries}\\
Global &
Anything true of every WeKnow kiosk everywhere. Rarely touched.\\
\rowcolor{weknowgreyL}
Country &
Things that follow the market --- currency, tax treatment, default
language.\\
Region (State) &
Regional variations within a country, where they exist.\\
\rowcolor{weknowgreyL}
City &
Things that follow geography --- relevant airports, default Uber
suggestions.\\
Location Group &
A set of sites treated the same, e.g.\ one hotel brand or one operator
across several properties.\\
\rowcolor{weknowgreyL}
Location &
The property itself --- hotel name, address, tax entity, banking details.\\
Kiosk Group &
A set of kiosks that share a specific configuration. \textbf{Not limited
to one location} --- a kiosk group can span sites, e.g.\ every Uber-only,
HEx-only or eSIM-only kiosk across the estate. Also used for kiosks
treated the same within a single site.\\
\rowcolor{weknowgreyL}
Kiosk &
Only what is unique to the single machine.\\
\end{tabularx}
\end{center}

The right-hand column is a guide to intent, not a fixed schema --- many
parameters can legitimately be set at more than one level. What matters is
choosing the widest one that is still correct.

\begin{wkinfo}[Check what a kiosk actually ends up with]
The kiosk's \textbf{Resolved} tab shows the effective configuration after
the whole hierarchy has been applied. Use it to confirm a kiosk is picking
up what you intended --- and to find the stale override when it is not.
\end{wkinfo}

%=====================================================================
% 3  PART 1 -- COMMAND CENTRE
%=====================================================================
\wkrolesection{Part 1 --- Command Centre}{Account Manager}

Done from any machine with a browser. You are creating the kiosk record,
clearing its prep checks, collecting the two credentials IT will need, and
scheduling the installation.

\subsection{Create the location and the kiosk}

\begin{enumerate}
  \item Open \href{https://command.weknow.technology}{Command Centre}
        and sign in. Ask the WKG Technology team if you do not have an
        account.

  \item Go to \href{https://command.weknow.technology/locations}{Locations}
        and create the location, working down the hierarchy:
        \textbf{Country $\rightarrow$ Region $\rightarrow$ City
        $\rightarrow$ Location}. Open it and \textbf{copy the Location
        ID} --- you need it when you create the installation record.
\end{enumerate}

\screenshot[0.50]{cc-location-id}{The location Overview tab. Use the copy icon next to \textbf{Location ID}.}

\begin{enumerate}\setcounter{enumi}{2}
  \item \emph{Only if required:} create a group from
        \href{https://command.weknow.technology/groups}{Groups}. Groups
        apply shared configuration to several sites or several kiosks at
        once:
    \begin{itemize}
      \item \textbf{Location groups} --- sites treated the same, e.g.\ one
            hotel brand or one operator across several properties.
      \item \textbf{Kiosk groups} --- kiosks that share a specific
            configuration. These are \emph{not} limited to a single
            location: a kiosk group can gather every Uber-only, HEx-only
            or eSIM-only kiosk across different sites.
    \end{itemize}
    Skip this for a single standalone kiosk.

  \item Go to \textbf{Kiosks} and select \textbf{Add Kiosk}:
    \begin{itemize}
      \item \textbf{Name} --- how the kiosk is referred to operationally,
            e.g.\ \emph{T2/T3 Tunnel - 2}.
      \item \textbf{Serial number} --- the workstation label,
            e.g.\ \wpath{WKG-POS-617}.
      \item \textbf{Hardware model} --- e.g.\ \emph{K10/20}.
      \item \textbf{Firmware version} --- e.g.\ \emph{1.0.0}.
    \end{itemize}
    Then select \textbf{Create Kiosk}.
\end{enumerate}

\screenshot[0.42]{cc-add-kiosk}{Add Kiosk. The kiosk is created unassigned --- it is tied to a location by the installation record.}

\begin{enumerate}\setcounter{enumi}{4}
  \item The new kiosk appears as \textbf{In Stock}. That is expected ---
        it stays In Stock until it is installed.

  \item If the kiosk belongs to any kiosk groups, open its \textbf{Edit}
        tab and assign them under \textbf{Group assignments}. A kiosk can
        be in several groups at once, and the order they sit in decides
        which one wins a conflict --- see \S\ref{sec:groupsorder}.
\end{enumerate}

\screenshot[0.36]{cc-group-assignments}{Group assignments on the kiosk Edit tab. Pick a group and select \textbf{Assign}; assigned groups are held in a display order, highest priority first.}

\subsection{Set the payment credentials}

Set the FreedomPay \textbf{Store ID (StID)} and \textbf{Terminal ID
(TID)} in the kiosk configuration, so payment requests are signed for the
correct merchant.

\subsection{Clear the prep checks}

Open the kiosk and go to the \textbf{Edit} tab. Mark the prep checks
\textbf{Complete} to confirm the kiosk is ready to config and install.

\screenshot[0.74]{cc-prep-checks}{Prep checks on the kiosk Edit tab.}

Before clearing them, verify with IT that:

\begin{itemize}
  \item all the hardware is plugged in and working --- printer, PDQ,
        sensors and scanners;
  \item FreedomPay has been installed, and the
        \wpath{FreewayClientService} and \wpath{FreewayServerService} are
        running.
\end{itemize}

\subsection{Collect the credentials for IT}

\begin{enumerate}
  \item On the kiosk's \textbf{Overview} tab, \textbf{copy the Kiosk ID}.
\end{enumerate}

\screenshot[0.48]{cc-kiosk-id}{Kiosk Overview --- copy the \textbf{Kiosk ID}.}

\begin{enumerate}\setcounter{enumi}{1}
  \item Switch to the \textbf{API keys} tab and select \textbf{Generate
        API key}. This is the key the kiosk uses to authenticate to the
        server.
\end{enumerate}

\screenshot[0.46]{cc-api-key}{Kiosk API keys. Generate the key here.}

\begin{wkwarn}[Copy the API key now]
The full key is displayed \textbf{only once, at creation}. Paste it
somewhere safe --- your password manager or the install ticket --- until
the installation is finished. If you lose it, generate a new one; the old
key cannot be recovered.
\end{wkwarn}

\begin{wkinfo}[Hand over to IT]
Send IT four things: the \textbf{Kiosk ID}, the \textbf{API key}, the
\textbf{Workstation ID} (\wpath{WKG-POS-XXX}), and the printer's COM port
if it is not the standard \wpath{COM8}. Send the API key by a private
channel, and delete it from that channel once the install is confirmed
working.
\end{wkinfo}

\subsection{Create the installation record}

The installation record ties the kiosk to its location and moves it out of
\emph{In Stock}.

\begin{enumerate}
  \item Go to
        \href{https://command.weknow.technology/installations}{Installations}
        and select \textbf{Add Installation}.

  \item Complete the form:
    \begin{itemize}
      \item \textbf{Kiosk} --- the kiosk you created above.
      \item \textbf{Location ID} --- the UUID copied earlier.
      \item \textbf{Planned date} --- the installation date.
      \item \textbf{Hardware model} and \textbf{Serial number} --- must
            match the kiosk record.
      \item \textbf{Installer name}, \textbf{Installer contact} and
            \textbf{Installation company}.
      \item \textbf{Notes} --- anything the engineer on site needs to
            know, e.g.\ \emph{``Install near main entrance''}.
    \end{itemize}

  \item Select \textbf{Create Installation}.
\end{enumerate}

\screenshot[0.36]{cc-add-installation}{Add Installation --- links the kiosk to its location and records who is installing it.}

\begin{wkinfo}[Done when]
The kiosk record exists with its prep checks cleared, the installation
record is created, and IT have the Kiosk ID, API key and Workstation ID.
\end{wkinfo}

%=====================================================================
% 3  PART 2 -- IT
%=====================================================================
\wkrolesection{Part 2 --- Install and lock the kiosk}{IT}

You need the Kiosk ID and API key from the Account Manager before you
start. The machine must already be built per the \emph{Kato Kiosk Setup
Guide}.

\subsection{Get the installer onto the machine}

\begin{enumerate}
  \item Download the latest \wpath{Kiosk Setup X.Y.Z.exe} from the
        \href{https://github.com/WeKnowTechnology/wkg-electron-app/releases}{app
        releases page}.

  \item Copy it into\\
        \wpath{C:\textbackslash Users\textbackslash SelfBooker\textbackslash Documents}\\
        This is the folder the N-central script looks in. Keep only one
        installer there --- if several are present the script takes the
        highest version and warns about the rest.
\end{enumerate}

Then install using \textbf{one} of the two methods below. The N-central
script is preferred: it validates the credentials, runs the install as
\wpath{SelfBooker} and checks the result for you.

\subsection{Option A --- One-run script from N-central}\label{sec:ncentral}

\begin{enumerate}
  \item Make sure \wpath{SelfBooker} is \textbf{signed in} on the machine.
        The install is per-user, and the script warns if the account is not
        logged on.

  \item Open the N-central remote connection to the machine. In the
        \textbf{System Shell} panel, select the \textbf{PowerShell} icon at
        the top right to switch the shell to PowerShell mode.
\end{enumerate}

\screenshot[0.94]{ncentral-powershell}{N-central System Shell. The PowerShell icon is the first of the three at the top right; the shell shows \textbf{Powershell} once switched.}

\begin{enumerate}\setcounter{enumi}{2}
  \item Edit the three values at the top of the script --- \wpath{\$KioskId},
        \wpath{\$WorkstationId} and \wpath{\$ApiKey} --- then paste the whole
        script into the shell.

  \item Press \textbf{Enter twice}: once to submit the script and once to
        run it.

  \item Wait for \textbf{\texttt{=== SUCCESS ===}}. Anything else ends in
        \texttt{=== FAILED ===} with the reason on the next line.
\end{enumerate}

\begin{lstlisting}
<#
    Deploy WKG Kiosk. Run as SYSTEM.
    Per-user install into the target account's profile via a scheduled task.
    Picks the highest-versioned installer present in the folder.
#>

$TargetUser     = 'SelfBooker'
$InstallerDir   = 'C:\Users\SelfBooker\Documents'
$KioskId        = '{KIOSK_ID_FROM_COMMAND_CENTRE}'
$WorkstationId  = '{WKG-POS-XXX}'
$ApiKey         = '{API_KEY_FROM_COMMAND_CENTRE}'
$RequireVersion = $null      # e.g. '1.1.1' to pin; $null accepts the newest found

$ErrorActionPreference = 'Stop'
$task  = 'KioskInstall'
$start = Get-Date

function Step($m) { Write-Host ("[{0,5:n1}s] {1}"       -f ((Get-Date)-$start).TotalSeconds, $m) }
function Ok($m)   { Write-Host ("[{0,5:n1}s]   OK  {1}" -f ((Get-Date)-$start).TotalSeconds, $m) -ForegroundColor Green }
function Warn($m) { Write-Host ("[{0,5:n1}s]   !   {1}" -f ((Get-Date)-$start).TotalSeconds, $m) -ForegroundColor Yellow }

try {
    Write-Host ""
    Write-Host "=== WKG Kiosk install: $WorkstationId on $env:COMPUTERNAME ==="
    Write-Host ""

    Step 'Validating parameters'
    $checks = @{
        KioskId       = @($KioskId,       '^[0-9a-fA-F]{8}-([0-9a-fA-F]{4}-){3}[0-9a-fA-F]{12}$')
        WorkstationId = @($WorkstationId, '^[A-Za-z0-9][A-Za-z0-9-]{1,14}$')
        ApiKey        = @($ApiKey,        '^kcs_[A-Za-z0-9_-]{24,}$')
    }
    foreach ($k in $checks.Keys) {
        $val, $pattern = $checks[$k]
        if ($val -cnotmatch $pattern) { throw "$k is malformed: '$val'" }
    }
    Ok 'Parameters well-formed'

    Step "Locating installer in $InstallerDir"
    if (-not (Test-Path -LiteralPath $InstallerDir)) { throw "Folder not found: $InstallerDir" }

    $candidates = Get-ChildItem -LiteralPath $InstallerDir -Filter 'Kiosk*Setup*.exe' -File -ErrorAction SilentlyContinue |
        ForEach-Object {
            if ($_.Name -match '(\d+\.\d+\.\d+)') {
                [pscustomobject]@{ Path = $_.FullName; Version = [version]$Matches[1]; Name = $_.Name }
            }
        } | Sort-Object Version -Descending

    if (-not $candidates) { throw "No Kiosk installer found in $InstallerDir" }

    $chosen        = @($candidates)[0]
    $InstallerPath = $chosen.Path
    Ok "Selected $($chosen.Name)  (v$($chosen.Version))"

    if (@($candidates).Count -gt 1) {
        $older = (@($candidates) | Select-Object -Skip 1 | ForEach-Object { $_.Name }) -join ', '
        Warn "$(@($candidates).Count) installers present. Ignoring: $older"
    }

    $fileVer = (Get-Item -LiteralPath $InstallerPath).VersionInfo.ProductVersion
    if ($fileVer -and ([version]$fileVer) -ne $chosen.Version) {
        Warn "Filename says v$($chosen.Version) but file metadata says v$fileVer"
    }
    if ($RequireVersion -and $chosen.Version -ne [version]$RequireVersion) {
        throw "Expected v$RequireVersion, found v$($chosen.Version)"
    }

    Step "Checking target user: $TargetUser"
    if (-not (Get-LocalUser -Name $TargetUser -ErrorAction SilentlyContinue)) {
        throw "No local user '$TargetUser'"
    }
    if ((quser 2>$null) -match "\s$TargetUser\s") { Ok 'User present and logged on' }
    else { Warn 'User exists but is not logged on - per-user install may misbehave' }

    Step 'Registering scheduled task'
    $argLine = "/S /KIOSKID=$KioskId /WORKSTATIONID=$WorkstationId /APIKEY=$ApiKey"
    Get-ScheduledTask -TaskName $task -ErrorAction SilentlyContinue | Unregister-ScheduledTask -Confirm:$false
    Register-ScheduledTask -TaskName $task `
        -Action    (New-ScheduledTaskAction -Execute $InstallerPath -Argument $argLine) `
        -Principal (New-ScheduledTaskPrincipal -UserId $TargetUser -LogonType Interactive -RunLevel Limited) `
        -Force | Start-ScheduledTask
    Ok 'Task registered and started'

    Step 'Waiting for installer to start (max 15s)'
    $by = (Get-Date).AddSeconds(15)
    do { Start-Sleep 1 }
    while ((Get-ScheduledTask -TaskName $task).State -ne 'Running' -and (Get-Date) -lt $by)

    $info = Get-ScheduledTask -TaskName $task | Get-ScheduledTaskInfo
    if ((Get-ScheduledTask -TaskName $task).State -ne 'Running' -and $info.LastTaskResult -ne 0) {
        throw ("Task did not start. LastTaskResult 0x{0:X}" -f $info.LastTaskResult)
    }
    Ok 'Installer running'

    Step 'Waiting for task to complete (max 120s)'
    $by = (Get-Date).AddSeconds(120)
    do { Start-Sleep 2 }
    while ((Get-ScheduledTask -TaskName $task).State -eq 'Running' -and (Get-Date) -lt $by)
    if ((Get-ScheduledTask -TaskName $task).State -eq 'Running') { Warn 'Task still running at timeout' }
    else { Ok 'Task finished' }

    Step 'Waiting for setup process to exit (max 60s)'
    $setupName = [System.IO.Path]::GetFileNameWithoutExtension($InstallerPath)
    $by = (Get-Date).AddSeconds(60)
    while ((Get-Process -Name $setupName -ErrorAction SilentlyContinue) -and (Get-Date) -lt $by) {
        Start-Sleep 2
    }
    if (Get-Process -Name $setupName -ErrorAction SilentlyContinue) {
        Warn 'Setup still alive (likely session 0, no desktop) - stopping it'
        Get-Process -Name $setupName -ErrorAction SilentlyContinue | Stop-Process -Force
    } else {
        Ok 'Setup exited cleanly'
    }
    Get-ScheduledTask -TaskName $task -ErrorAction SilentlyContinue | Unregister-ScheduledTask -Confirm:$false

    Step 'Verifying install'
    $exe  = "C:\Users\$TargetUser\AppData\Local\Programs\Kiosk\Kiosk.exe"
    $json = "C:\Users\$TargetUser\AppData\Roaming\WKG Kiosk\kiosk-identity.json"

    if (-not (Test-Path -LiteralPath $exe)) { throw "Kiosk.exe missing at $exe" }
    Ok ("Kiosk.exe present ({0} MB)" -f [math]::Round((Get-Item -LiteralPath $exe).Length / 1MB))

    if (-not (Test-Path -LiteralPath $json)) { throw "Not provisioned, no file at $json" }
    Ok 'Identity file written'

    $cfg = Get-Content -LiteralPath $json -Raw | ConvertFrom-Json
    $bad = @()
    if ($cfg.kioskId       -cne $KioskId)       { $bad += "kioskId='$($cfg.kioskId)'" }
    if ($cfg.workstationId -cne $WorkstationId) { $bad += "workstationId='$($cfg.workstationId)'" }
    if ($bad) { throw "Config mismatch: $($bad -join '; ')" }

    Ok "Identity matches: $WorkstationId"
    Ok "Printer: $($cfg.printerDestination)"

    Write-Host ""
    Write-Host ("=== SUCCESS in {0:n1}s  (v{1}) ===" -f ((Get-Date)-$start).TotalSeconds, $chosen.Version) -ForegroundColor Green
    Write-Host ""
    exit 0
}
catch {
    Write-Host ""
    Write-Host ("=== FAILED after {0:n1}s ===" -f ((Get-Date)-$start).TotalSeconds) -ForegroundColor Red
    Write-Host $_.Exception.Message -ForegroundColor Red
    Write-Host ""
    Get-ScheduledTask -TaskName $task -ErrorAction SilentlyContinue | Unregister-ScheduledTask -Confirm:$false
    exit 1
}
\end{lstlisting}

\begin{wkinfo}[The script does not set the printer]
It installs with the default port, \wpath{COM8}, and prints the port it
ended up with on the \texttt{Printer:} line. If the printer is elsewhere,
correct it in \wpath{kiosk-identity.json} --- see \S\ref{sec:comport}.
\end{wkinfo}

Skip to step~3 of Option~B to confirm the identity file, then continue with
the lockdown.

\subsection{Option B --- Manually, signed in as \wpath{SelfBooker}}

Done \textbf{signed in as \wpath{SelfBooker}} --- the account the kiosk
runs as.

\begin{enumerate}
  \item Sign in to Windows as \wpath{SelfBooker}.

  \item Open \textbf{PowerShell} in
        \wpath{C:\textbackslash Users\textbackslash SelfBooker\textbackslash Documents}
        and run the silent install, substituting the installer version and
        your four values:
\end{enumerate}

\begin{lstlisting}
& ".\Kiosk Setup X.Y.Z.exe" /S `
  /KIOSKID={KIOSK_ID_FROM_COMMAND_CENTRE} `
  /WORKSTATIONID={WKG-POS-XXX} `
  /APIKEY={API_KEY_FROM_COMMAND_CENTRE} `
  /PRINTER=COM8
\end{lstlisting}

\begin{center}
\renewcommand{\arraystretch}{1.3}\small
\begin{tabularx}{\textwidth}{L{34mm} X}
\rowcolor{weknownight}\color{white}\textbf{Switch} &
\color{white}\textbf{Value}\\
\wpath{/S} & Silent install --- no prompts. Always used.\\
\rowcolor{weknowgreyL}
\wpath{/KIOSKID} & The Kiosk ID from the Account Manager.\\
\wpath{/WORKSTATIONID} & The machine label, e.g.\ \wpath{WKG-POS-617}.\\
\rowcolor{weknowgreyL}
\wpath{/APIKEY} & The API key from the Account Manager.\\
\wpath{/PRINTER} & The printer's COM port. \wpath{COM8} on a standard
build. If nothing prints, the port is wrong --- see
\S\ref{sec:comport}.\\
\end{tabularx}
\end{center}

\begin{enumerate}\setcounter{enumi}{2}
  \item Confirm the app wrote its identity file. Open
        \wpath{kiosk-identity.json} (\S\ref{sec:paths}) and run through the
        checks in \S\ref{sec:identity}. If the file is missing, create it
        before continuing --- the app will not connect without it.

  \item Sign out of \wpath{SelfBooker}, then sign in as \wpath{WKGAdmin}
        for the lockdown below.
\end{enumerate}

\subsection{Lock the machine into kiosk mode}

The Shell Launcher utility replaces the Windows desktop with the Kato app
for the \wpath{SelfBooker} account, so the machine boots straight into the
kiosk and there is no desktop to escape to.

\begin{enumerate}
  \item Signed in as \wpath{WKGAdmin}, run\\
        \wpath{C:\textbackslash Kiosk\textbackslash SSM25\_files\textbackslash Shell-Launcher-Utilities\textbackslash SLUtility.exe}

  \item Select \textbf{Disable User Shell}.

  \item \textbf{Close and reopen} \wpath{SLUtility}. It needs the restart
        to pick up the change --- do not skip this.

  \item Select the \wpath{SelfBooker} user and change the shell path to
        the app, leaving the default \textbf{Action 3}:\\
        \wpath{C:\textbackslash Users\textbackslash SelfBooker\textbackslash AppData\textbackslash Local\textbackslash Programs\textbackslash Kiosk\textbackslash Kiosk.exe}

  \item Select \textbf{Set User Shell}.

  \item Select \textbf{Enable Shell Launcher}.

  \item Sign out of \wpath{WKGAdmin} and sign in as \wpath{SelfBooker}.
        The Kato app should launch full screen with no desktop behind it.
        Reboot once to confirm it comes back the same way.
\end{enumerate}

\begin{wkwarn}[Order matters]
Disable $\rightarrow$ restart the utility $\rightarrow$ set the path
$\rightarrow$ set shell $\rightarrow$ enable. Do not set the shell before
restarting \wpath{SLUtility}.
\end{wkwarn}

\begin{wkinfo}[Done when]
The machine reboots straight into the Kato app as \wpath{SelfBooker}, with
no visible desktop, taskbar or file explorer.
\end{wkinfo}

%=====================================================================
% 4  FINAL CHECKLIST
%=====================================================================
\wksection{Final checklist}

Run through this before leaving site.

\begin{center}
\renewcommand{\arraystretch}{1.45}\small
\begin{tabularx}{\textwidth}{L{8mm} L{14mm} X}
\rowcolor{weknownight}\color{white} & \color{white}\textbf{Who} &
\color{white}\textbf{Check}\\
\emptybox & AM & Kiosk record has the correct name, serial and hardware
model.\\
\rowcolor{weknowgreyL}
\emptybox & AM & All prep checks (hardware, firmware, software) cleared.\\
\emptybox & AM & Installation record created and linked to the location;
kiosk no longer \emph{In Stock}.\\
\rowcolor{weknowgreyL}
\emptybox & AM & API key deleted from the ticket / chat once the install
is confirmed working.\\
\emptybox & IT & \wpath{kiosk-identity.json} present and correct.\\
\rowcolor{weknowgreyL}
\emptybox & IT & Machine reboots straight into the Kato app --- no
desktop.\\
\emptybox & IT & Printer prints, scanner scans, sensors respond.\\
\rowcolor{weknowgreyL}
\emptybox & IT & Test payment taken on the PDQ and refunded.\\
\end{tabularx}
\end{center}

%=====================================================================
% 5  TROUBLESHOOTING
%=====================================================================
\wksection{When things go wrong}

\begin{center}
\renewcommand{\arraystretch}{1.35}\small
\begin{tabularx}{\textwidth}{L{50mm} X}
\rowcolor{weknownight}\color{white}\textbf{Symptom} &
\color{white}\textbf{First thing to try}\\
Machine boots to the Windows desktop, not the app &
The \wpath{SLUtility} sequence was not followed. Redo it in order,
restarting the utility between \emph{Disable User Shell} and setting the
path.\\
\rowcolor{weknowgreyL}
App launches but will not connect &
Work through the \wpath{kiosk-identity.json} checks in
\S\ref{sec:identity} --- dev server URLs, a stray backtick on a value, a
wrong Kiosk ID, or an API key regenerated after the install.\\
Installer runs but nothing appears &
The switches are case-sensitive and the backtick line continuations must
be the last character on each line. Retype the command rather than
pasting it from a chat client, or use the N-central script
(\S\ref{sec:ncentral}).\\
N-central script ends in \texttt{FAILED} &
The line under \texttt{FAILED} names the cause. \emph{Malformed} means a
value at the top of the script is mistyped; \emph{No Kiosk installer
found} means the installer is not in
\wpath{C:\textbackslash Users\textbackslash SelfBooker\textbackslash Documents};
\emph{Not provisioned} usually means \wpath{SelfBooker} was not signed in.\\
\rowcolor{weknowgreyL}
Printer does nothing &
The configured COM port is wrong. Find the port the printer is actually
on and correct it in \wpath{kiosk-identity.json} --- no reinstall needed
(\S\ref{sec:comport}).\\
Prep checks will not clear &
They are marked manually on the kiosk \textbf{Edit} tab --- they do not
clear themselves when the underlying condition is met.\\
\rowcolor{weknowgreyL}
Payment fails on the test transaction &
That is a machine-build issue, not a configuration one. See the
\emph{Kato Kiosk Setup Guide}.\\
\end{tabularx}
\end{center}

\subsection{Checking \wpath{kiosk-identity.json}}\label{sec:identity}

Most \emph{will not connect} faults are in this file. Open it in Notepad
(path in \S\ref{sec:paths}). A correct file looks like this:

\begin{lstlisting}
{
  "kioskId": "{KIOSK_ID_FROM_COMMAND_CENTRE}",
  "apiBaseUrl": "https://api.weknow.technology/api",
  "apiKey": "{API_KEY_FROM_COMMAND_CENTRE}",
  "workstationId": "WKG-POS-XXX",
  "printerDestination": "COM8",
  "fleetSocketUrl": "wss://api.weknow.technology"
}
\end{lstlisting}

\begin{center}
\renewcommand{\arraystretch}{1.35}\small
\begin{tabularx}{\textwidth}{L{34mm} X}
\rowcolor{weknownight}\color{white}\textbf{Field} &
\color{white}\textbf{Check}\\
\wpath{apiBaseUrl} &
Must be \wpath{https://api.weknow.technology/api} --- \textbf{not}
\wpath{https://dev.weknow.technology/api}.\\
\rowcolor{weknowgreyL}
\wpath{fleetSocketUrl} &
Must be \wpath{wss://api.weknow.technology} --- \textbf{not}
\wpath{wss://dev.weknow.technology}.\\
\wpath{kioskId} &
Matches the Kiosk ID on the kiosk's Overview tab, and does \textbf{not end
in a backtick} (\wpath{`}).\\
\rowcolor{weknowgreyL}
\wpath{apiKey} &
Matches the key generated for this kiosk, and does \textbf{not end in a
backtick}.\\
\wpath{workstationId} &
The machine label, e.g.\ \wpath{WKG-POS-617} --- not the Kiosk ID.\\
\rowcolor{weknowgreyL}
\wpath{printerDestination} &
The printer's COM port (\S\ref{sec:comport}).\\
\end{tabularx}
\end{center}

\begin{wkwarn}[Why a backtick ends up in a value]
In the manual install command, each line ends in a backtick to continue
it. If a space follows the backtick, or the command is pasted onto one
line, PowerShell passes the backtick through as part of the value. Delete
it from the file.
\end{wkwarn}

Fix any wrong value directly in the file, save, and restart the app ---
there is no need to reinstall. The one exception is an API key that no
longer works at all: generate a new one in Command Centre and paste that in.

\subsection{Finding the printer's COM port}\label{sec:comport}

\begin{wkwarn}[Device Manager is not the answer here]
The app writes \textbf{straight to the COM port} --- there is no printer
driver in between. The port the printer actually answers on does
\emph{not} always match what Device Manager shows against it, so read the
ports, then test them.
\end{wkwarn}

\paragraph{1. List the ports Windows can see.} In PowerShell:

\begin{lstlisting}
# Just the port names
[System.IO.Ports.SerialPort]::GetPortNames()

# ...with the device name against each, which narrows the field
Get-CimInstance Win32_PnPEntity |
  Where-Object { $_.Name -match 'COM\d+' } |
  Select-Object Name, DeviceID
\end{lstlisting}

Or, from Command Prompt, \wpath{mode} lists every COM device with its
current settings:

\begin{lstlisting}
mode
\end{lstlisting}

\paragraph{2. Test each port until the printer reacts.} Writing anything
to the right port makes the printer spit out a few lines of nonsense ---
that is the confirmation. Run this from \textbf{Command Prompt}, one port
at a time, working through the list from step 1:

\begin{lstlisting}
DIR > COM3
DIR > COM4
DIR > COM8
\end{lstlisting}

From PowerShell, hand the redirect to \wpath{cmd} instead, so it reaches
the device rather than creating a file:

\begin{lstlisting}
cmd /c "DIR > COM8"
\end{lstlisting}

How to read the result:

\begin{center}
\renewcommand{\arraystretch}{1.3}\small
\begin{tabularx}{\textwidth}{L{54mm} X}
\rowcolor{weknownight}\color{white}\textbf{What happens} &
\color{white}\textbf{Means}\\
Printer prints garbage &
That is the port. Note it and move to step 3.\\
\rowcolor{weknowgreyL}
Command returns silently, nothing prints &
The port exists but the printer is not on it. Try the next one.\\
\emph{Access is denied} &
Something already holds the port --- usually the kiosk app. Close it and
try again.\\
\rowcolor{weknowgreyL}
\emph{The system cannot find the path specified} &
No such port. Try the next one.\\
\end{tabularx}
\end{center}

\paragraph{3. Set the port in \wpath{kiosk-identity.json}.} There is
\textbf{no need to rerun the silent install}. Open
\wpath{kiosk-identity.json}, set the printer port to the one you
confirmed, save, then restart the app and print a test receipt.

\begin{wkinfo}[Note it on the installation record]
If the printer turned out to be on a non-standard port, put it in the
\textbf{Notes} field of the installation record. It saves the next
engineer repeating this.
\end{wkinfo}

%=====================================================================
% APPENDIX -- ONE-PAGE QUICK REFERENCE
%=====================================================================
\clearpage
\appendix
\wksection{Quick reference}

One page for someone who has done this before. Work down your column;
the full detail is in the sections above.

\begin{center}
\begin{minipage}[t]{0.475\textwidth}
\colorbox{weknowsky}{\parbox{\dimexpr\textwidth-2\fboxsep}{%
  \vspace{2pt}\centering\bfseries\color{weknownight}
  Part 1 --- Account Manager\vspace{2pt}}}
\vspace{4pt}
\footnotesize
\begin{enumerate}[leftmargin=6mm,itemsep=3pt]
  \item \emptybox\ Collect the location information (\S1.1). Skip RPS
        customer code and outlet code --- legacy.
  \item \emptybox\ Create the location; copy the \textbf{Location ID}.
  \item \emptybox\ Create groups if needed. Kiosk groups can span sites.
  \item \emptybox\ \textbf{Add Kiosk} --- name, serial, hardware model,
        firmware.
  \item \emptybox\ Assign kiosk groups on the \textbf{Edit} tab.
        Top of the list wins a conflict.
  \item \emptybox\ Set the \textbf{Store ID} and \textbf{Terminal ID}.
  \item \emptybox\ Clear the prep checks --- only after IT confirm
        hardware and FreedomPay.
  \item \emptybox\ Copy the \textbf{Kiosk ID} (Overview tab).
  \item \emptybox\ Generate the \textbf{API key} --- shown once. Copy it
        now.
  \item \emptybox\ Send the four values to IT (below).
  \item \emptybox\ Create the \textbf{installation record}.
\end{enumerate}
\end{minipage}\hfill
\begin{minipage}[t]{0.475\textwidth}
\colorbox{weknowgreyL}{\parbox{\dimexpr\textwidth-2\fboxsep}{%
  \vspace{2pt}\centering\bfseries\color{weknownight}
  Part 2 --- IT\vspace{2pt}}}
\vspace{4pt}
\footnotesize
\begin{enumerate}[leftmargin=6mm,itemsep=3pt]
  \item \emptybox\ Confirm the machine is built --- see the
        \emph{Kato Kiosk Setup Guide}.
  \item \emptybox\ Sign in to Windows as \wpath{SelfBooker}.
  \item \emptybox\ Copy the latest installer into
        \wpath{SelfBooker\textbackslash Documents}.
  \item \emptybox\ Run the N-central script, or the silent install below.
  \item \emptybox\ Check \wpath{kiosk-identity.json} --- IDs, API key, no
        trailing backtick, \wpath{api.} not \wpath{dev.} URLs.
  \item \emptybox\ Sign out; sign in as \wpath{WKGAdmin}.
  \item \emptybox\ \wpath{SLUtility}: \textbf{Disable User Shell}.
  \item \emptybox\ Close and reopen \wpath{SLUtility}. Do not skip.
  \item \emptybox\ Select \wpath{SelfBooker}, set the shell path,
        keep Action 3.
  \item \emptybox\ \textbf{Set User Shell}, then \textbf{Enable Shell
        Launcher}.
  \item \emptybox\ Sign in as \wpath{SelfBooker}; reboot and confirm the
        app comes back.
  \item \emptybox\ Test print, test scan, test payment.
\end{enumerate}
\end{minipage}
\end{center}

\vspace{6pt}

\begin{center}
\renewcommand{\arraystretch}{1.2}\footnotesize
\begin{tabularx}{\textwidth}{L{34mm} X}
\rowcolor{weknownight}\color{white}\textbf{AM hands IT} &
\color{white}\textbf{What it is}\\
Kiosk ID & UUID from the kiosk Overview tab.\\
\rowcolor{weknowgreyL}
API key & Generated on the API keys tab. Shown once.\\
Workstation ID & The machine label, \wpath{WKG-POS-XXX}.\\
\rowcolor{weknowgreyL}
Printer COM port & \wpath{COM8} unless the site differs.\\
\end{tabularx}
\end{center}

\vspace{4pt}
\noindent\textbf{\footnotesize Silent install} --- PowerShell, as
\wpath{SelfBooker}:

\begin{lstlisting}
& ".\Kiosk Setup X.Y.Z.exe" /S /KIOSKID={KIOSK_ID} /WORKSTATIONID={WKG-POS-XXX} /APIKEY={API_KEY} /PRINTER=COM8
\end{lstlisting}

\begin{center}
\renewcommand{\arraystretch}{1.2}\footnotesize
\begin{tabularx}{\textwidth}{L{34mm} X}
\rowcolor{weknownight}\color{white}\textbf{Path} &
\color{white}\textbf{Where}\\
Kato app &
\wpath{C:\textbackslash Users\textbackslash SelfBooker\textbackslash AppData\textbackslash Local\textbackslash Programs\textbackslash Kiosk\textbackslash Kiosk.exe}\\
\rowcolor{weknowgreyL}
Shell Launcher utility &
\wpath{C:\textbackslash Kiosk\textbackslash SSM25\_files\textbackslash Shell-Launcher-Utilities\textbackslash SLUtility.exe}\\
Find the printer port &
List with \wpath{mode}, then \wpath{DIR > COM3} from Command Prompt.
Fix in \wpath{kiosk-identity.json} --- no reinstall.\\
\end{tabularx}
\end{center}

\end{document}
